% ==================================================================
% id:       seq-031
% topic:    סדרות
% subtopic: אריתמטיקה של גבולות — הוכח או הפרך
% solution: full
% ==================================================================
\textbf{תהיינה
$(a_n)_{n=1}^{\infty}$
ו-$(b_n)_{n=1}^{\infty}$
סדרות.
הוכיחו או הפריכו:}
\begin{enumerate}
\item
אם
$a_n\to L$
ו-$a_n\ne0$
לכל
$n$,
אז
$L\ne0$.
\hint{האם סדרה שכל איבריה שונים מאפס יכולה לשאוף לאפס?}
\begin{solbox}
הטענה לא נכונה. הסדרה
$a_n=\frac1n$
מקיימת
$a_n\ne0$
לכל
$n$,
אבל
$a_n\to0$.
\end{solbox}

\item
אם
$a_n+b_n\to0$
ו-$a_n-b_n\to0$,
אז
$a_n\to0$
ו-$b_n\to0$.
\hint{נסו להשתמש באריתמטיקה, והיעזרו בסדרות שידוע שהן מתכנסות.}
\begin{solbox}
הטענה נכונה. לכל
$n$
מתקיים
$$a_n=\frac12\bigl[(a_n+b_n)+(a_n-b_n)\bigr],\qquad b_n=\frac12\bigl[(a_n+b_n)-(a_n-b_n)\bigr].$$
לפי אריתמטיקה של גבולות (סכום והפרש של סדרות מתכנסות, וכפל בקבוע),
$$a_n\goes\frac12(0+0)=0,\qquad b_n\goes\frac12(0-0)=0 .$$
\end{solbox}

\item
אם
$(a_n)$
מתכנסת, אז
$a_{n+1}-a_n\to0$.
\hint{אם $a_n\to L$, לאן שואפת $a_{n+1}$? היעזרו בשאלה \qref{seq-030}.}
\begin{solbox}
הטענה נכונה. נניח ש-$a_n\to L$.
לפי סעיף ג' של שאלה
\qref{seq-030},
גם
$a_{n+1}\to L$.
לפי אריתמטיקה של גבולות (הפרש של סדרות מתכנסות),
$$a_{n+1}-a_n\goes L-L=0 .$$
\end{solbox}

\item
אם
$a_{n+1}-a_n\to0$,
אז
$(a_n)$
מתכנסת.
\hint{נסו לחשוב על חזקות של $n$.}
\begin{solbox}
הטענה לא נכונה. ניקח
$a_n=\sqrt n$.

\textbf{ההפרשים שואפים ל-$0$:}
בעזרת כפל בצמוד,
$$0<\sqrt{n+1}-\sqrt n=\frac{1}{\sqrt{n+1}+\sqrt n}<\frac{1}{\sqrt n}=\lr{\frac1n}^{1/2}.$$
מכיוון ש-$\frac1n\to0$, לפי אריתמטיקה של חזקות
$\lr{\frac1n}^{1/2}\goes0^{1/2}=0$.
לפי כלל הסנדוויץ', $a_{n+1}-a_n=\sqrt{n+1}-\sqrt n\goes0$.

\textbf{הסדרה אינה מתכנסת:}
היא אינה חסומה, כי לכל
$M>0$
ולכל
$n>M^2$
מתקיים
$\sqrt n>M$.
סדרה מתכנסת היא חסומה, ולכן
$\sqrt n$
אינה מתכנסת.
\end{solbox}

\end{enumerate}
